\chapter{בניית הצוות המייסד}

המשאב הקריטי ביותר בשלב המוקדם אינו ההון אלא הצוות. פיטר תיל מדגיש שראיון
המפתח אינו ``מדוע הרעיון טוב'', אלא ``מדוע דווקא הצוות הזה'' — האם המייסדים
מכירים זה את זה היטב, משלימים זה את זה, ומחויבים לטווח ארוך~\cite{thiel2014zero}.

\section{תפקידים משלימים}

צוות מייסד בריא משלב יכולות בנייה (\en{building}) ויכולות הפצה (\en{distribution}).
חברה רבת-עוצמה מבחינה טכנולוגית אך חסרת ערוץ הפצה תיכשל לא פחות מחברה עם שיווק
מצוין ומוצר חלש. חלוקת התפקידים צריכה להיות ברורה, אך גבולות הגזרה גמישים דיים
כדי לאפשר עבודה משותפת בתנאי אי-ודאות.

\section{הון אנושי, מניות והסכמות מוקדמות}

חלוקת המניות בין המייסדים והגדרת תקופת הבשלה \term{(vesting)} הן החלטות שכדאי
לקבל מוקדם ובכתב, בעוד היחסים טובים. הסכמות לא-פורמליות הן מקור נפוץ לסכסוכים
מאוחרים שעלולים לפרק את החברה. עיקרון מנחה: לתגמל מחויבות עתידית ולא רק תרומה
היסטורית, וליישר את התמריצים של כל המייסדים עם הצלחת החברה לטווח ארוך.

\section{תרבות וגיוסי העובדים הראשונים}

לאחר המייסדים, עשרת העובדים הראשונים מעצבים את \emph{התרבות} של החברה לשנים
הבאות. בשלב זה כל גיוס הוא החלטה אסטרטגית: עובד לא-מתאים פוגע בצוות קטן הרבה
יותר מאשר בארגון גדול. עדיף לגייס לאט, להעדיף יכולת למידה ותחושת בעלות
\term{(ownership)} על פני ניסיון צר, ולחפש אנשים המתפקדים היטב באי-ודאות.
תרבות אינה סיסמאות על הקיר אלא סך ההתנהגויות שהמייסדים מתגמלים ומענישים בפועל;
היא נקבעת מוקדם, וקשה מאוד לשנותה מאוחר.

\section{קצב, מיקוד ומשמעת ביצוע}

בשלב המוקדם היתרון התחרותי הגדול ביותר של סטארט-אפ הוא \emph{קצב} — היכולת
ללמוד ולשחרר מהר יותר מהמתחרים. קצב נשמר באמצעות מיקוד אכזרי: לבחור בכל רגע
את המשימה האחת החשובה ביותר, ולסרב לעשרות הזדמנויות סבירות שמפזרות את הצוות.
פגישות שבועיות קצרות סביב מדד-על אחד, יחד עם הגדרת ``הצעד הבא הקטן ביותר''
שניתן לשחרר, ממירות חזון מעורפל לתנועה מדידה קדימה.
